\documentclass[10pt,a4paper]{amsart}
\usepackage[margin=19mm]{geometry}
\usepackage{amsmath,amssymb,mathtools}
\usepackage[scr=rsfs,cal=euler]{mathalfa}
\usepackage{xfrac}
\usepackage[unicode,hidelinks,bookmarks=false]{hyperref}
\newcommand{\ie}{i.e.\ }
\newcommand{\cf}{cf.\ }
\newcommand{\resp}{respectively\ }
\newcommand{\Subsection}{\subsection}
\providecommand{\nobreakdash}{\nobreak\hbox{-}}
\setlength{\emergencystretch}{2em}
\multlinegap0pt
\usepackage{amssymb}
\usepackage{xcolor}
\usepackage{algorithm}\usepackage{algpseudocode}
\usepackage{tikz}
\newtheorem{remark}{Remark}
\newtheorem{theorem}{Theorem}
\title{$F$-Contraction with an Auxiliary Function and Its Application to Terrain-Following Airplane Navigation}
\author{Irom Shashikanta Singh, Yumnam Mahendra Singh}
\date{Source: arXiv:2603.10556v1 (CC BY 4.0)}
\begin{document}
\maketitle

\noindent\emph{Source and licence.} $F$-Contraction with an Auxiliary Function and Its Application to Terrain-Following Airplane Navigation. Version arXiv:2603.10556v1 (CC BY 4.0), distributed by arXiv under the Creative Commons Attribution 4.0 International licence, http://creativecommons.org/licenses/by/4.0/, as recorded in the arXiv licence metadata for this exact version and shown on the abstract page https://arxiv.org/abs/2603.10556v1. Full licence notice: \texttt{licenses/arxiv-2603.10556-CC-BY-4.0.txt}.\par\smallskip

\noindent\textit{Editorial scope.} Counted unit: the article's theorem appth (source0.tex lines 1591--1600). The article's complete original proof of it is delivered with it (source0.tex lines 1601--1711, one \texttt{proof} environment of 111 lines). No mathematics is written, repaired or re-derived: the statement and the proof are the article's own lines, in the article's own notation, and no retained line is substituted or re-flowed.  Retained uncounted as the source-local context the proof needs: lines 386--394; lines 598--606; lines 1471--1474; lines 1500--1508.  Nothing was omitted from any retained span, so the package carries no omission record.  No retained line contains a citation, so the wrapper carries no bibliography and no bibliography entry.  No retained line contains a figure environment, an image-inclusion command or drawing code, so no image is delivered and the package declares no asset dependency.  Editorial formatting only, in addition: an AMS \texttt{amsart} preamble that restates the article's own declarations and macros used by the retained lines, namely the environments \texttt{remark}, \texttt{theorem} on separate counters, together with the standard packages those lines need.  The retained environments print in this document as Remark 1, Theorem 1 in order of appearance, whereas the article prints the same items with the numbering of its own full document.  The complete native source of the article as distributed in the arXiv e-print for this exact version is bundled unchanged as \texttt{sources/196145/source0.tex}.
\begin{remark}\label{remark1} Under different mapping $F \in \mathcal{F}$, we obtain the following consequences:
    \begin{enumerate}
        \item [1.] If $F(\xi)=\ln \xi $, \textit{$S^F$-contraction} reduces to \textit{$S^B$-contraction} with $\beta=e^{-\omega  }$.
        \item[2.]  If $F(\xi)=\ln \xi  +\xi$, \textit{$S^F$-contraction} reduces to
        \begin{equation*}
            \frac{\eth(\Upsilon \xi,S(\Upsilon \xi,\Upsilon \zeta))+\eth(\Upsilon \zeta,S(\Upsilon \xi,\Upsilon \zeta))}{\eth(\xi,S(\xi,\zeta))+\eth(\zeta ,S(\xi,\zeta))}\cdot\frac{e^{\eth(\Upsilon \xi,S(\Upsilon \xi,\Upsilon \zeta))+\eth(\Upsilon \zeta,S(\Upsilon \xi,\Upsilon \zeta))}}{e^{\eth(\xi,S(\xi,\zeta))+\eth(\zeta ,S(\xi,\zeta))}}\le e^{-\omega  }.
        \end{equation*}
    \end{enumerate}
\end{remark}

\begin{theorem}\label{sf theorem}
  Let $(W,\eth,s)$ be a complete super-metric space and $S:W\times W \to W$ be a given mapping. 
If $\Upsilon  : W \to W$ is an \textit{$S^F$-contraction}, then $\Upsilon $ is a \textit{Picard operator} provided that, for all $u,v \in W$,
\[
\lim_{n \to \infty} \eth(\Upsilon^n u, v) = 0 
\quad \Longrightarrow \quad 
\lim_{n \to \infty} \eth(\Upsilon(\Upsilon^n u), \Upsilon  v) = 0.
\]
\end{theorem}

\begin{equation}\label{aa2}
\varkappa_{n+1}(\xi)
= -\gamma(\xi)+(\mathcal{G}K+I)\varkappa_n(\xi),
\end{equation}

\begin{equation}\label{aa4}
\frac{d^3\sigma}{dt^3}
+a_2\frac{d^2\sigma}{dt^2}
+a_1\frac{d\sigma}{dt}
=
b_2\frac{d^2\varkappa}{dt^2}
+b_1\frac{d\varkappa}{dt}
-b_0\varkappa.
\end{equation}

\begin{theorem}\label{appth}
     Suppose that there exists $\omega  ,\delta,a>0$ and $p\ge 1$ and a sequence $\Delta_n=|\varkappa_{n}-\varkappa_{n-1}-a|^p$ such that the conditions given below are satisfied.

\begin{itemize}
\item[(F$_1$)] $(\xi - \xi_0)S\delta \le 1 + e^{\frac{1}{p}(\Delta_n-\Delta_{n+1}-\omega  )},$
\item[(F$_2$)] $\dfrac{\mathcal{G}K\varkappa_{n}(\xi) - \mathcal{G}K\varkappa_{n-1}(\xi)}{\varkappa_n - \varkappa_{n-1}-a} < -1 + e^{\frac{1}{p}(\Delta_n-\Delta_{n+1}-\omega  )}.$
\end{itemize}

Then, \eqref{aa2} converges.
\end{theorem}

\begin{proof}
    
 Let us define the metric $\eth(f(\xi), g(\xi)) = \max |f(\xi) - g(\xi)|^p,p\ge1$, which is a complete super-metric space with $s=2^{p-1}$.

If the operator on the right-hand side of \eqref{aa2} is an $S^F$-contraction, then it possesses a unique fixed point $\xi^*$ given by
\[
\xi^* = -\gamma + (\mathcal{G}K + I)\xi^*.
\]

Taking $S(\xi,\zeta)=\zeta+ a, a> 0$, we have
\begin{equation}\label{a1}
\begin{split}
|\varkappa_{n+1}(\xi)-S(\varkappa_{n+1}(\xi),&\varkappa_{n}(\xi))|^p+|\varkappa_{n}(\xi)-S(\varkappa_{n+1}(\xi),\varkappa_{n}(\xi))|^p\\
&= \Delta_{n+1}+a^p\\
&= \left|\varkappa_{n+1}(\xi) - \varkappa_{n}(\xi)-a\right|^p+a^p\\
&= \left|(\mathcal{G}K+I)\varkappa_{n}(\xi) - (\mathcal{G}K+I)\varkappa_{n-1}(\xi)-a\right|^p+a^p  \\
&= \left|\mathcal{G}K\varkappa_n(\xi) - \mathcal{G}K\varkappa_{n-1}(\xi) + \varkappa_n(\xi) - \varkappa_{n-1}(\xi)-a\right|^p+a^p \\
&= \left| \frac{\mathcal{G}K\varkappa_n(\xi) - \mathcal{G}K\varkappa_{n-1}(\xi)}{\varkappa_n(\xi) - \varkappa_{n-1}(\xi)-a} +1 \right|^p 
\left|\varkappa_n(\xi) - \varkappa_{n-1}(\xi)-a\right|^p+a^p. 
\end{split}
\end{equation}


Consider
\[
\frac{\eth(\tan \sigma)}{d\varkappa} = \sec^2\sigma\frac{d\sigma}{d\varkappa} = (1 + \tan^2 \sigma)\frac{d\sigma}{dt}\frac{dt}{d\varkappa}.
\]

This can be rewritten as:
\[
\left(1 + \left(\frac{dh}{d\xi} \right)^2\right)
\left(\frac{d\sigma}{dt}\right)\left(\frac{dt}{d\varkappa}\right).
\]

Hence, the quantity $\frac{dh}{d\xi}$ remains bounded due to the restrictions on the aircraft's acceleration, while $\frac{d\sigma}{dt}$ and $\frac{d\varkappa}{dt}$ are related through Eq.~\eqref{aa4}, representing the governing dynamics of the aircraft.

Now, we will show that \[
\frac{\mathcal{G}K\varkappa_n(\xi) - \mathcal{G}K\varkappa_{n-1}(\xi)}{\varkappa_n - \varkappa_{n-1}-a} > -1 - e^{\frac{1}{p}(\Delta_n-\Delta_{n+1}-\omega  )}.
\]
From condition (F$_1$), we have
\begin{equation}\label{a3}
\begin{split}
 &(\xi - \xi_0)S\delta \le 1 + e^{\frac{1}{p}(\Delta_n-\Delta_{n+1}-\omega  )} \\
&\Rightarrow  -(\xi - \xi_0)S\delta >- 1 - e^{\frac{1}{p}(\Delta_n-\Delta_{n+1}-\omega  )}. 
\end{split}
\end{equation}

Let
\begin{equation*}
\frac{\eth(\tan \sigma)}{d\varkappa} \ge -\delta , 
\end{equation*}
where $\delta$ is a positive constant, which has not been specified yet.

Now, consider the difference between $\tan \sigma_n$ and $\tan (\sigma_{n-1}+a)$:
\[
\tan \sigma_n - \tan (\sigma_{n-1}+a) = \int_{\varkappa_{n-1}+a}^{\varkappa_n} \frac{\eth(\tan \sigma)}{d\varkappa} d\varkappa \ge -\eth(\varkappa_n-\varkappa_{n-1}-a).
\]

Therefore, we have:
\begin{equation}\label{a4}
\frac{S \int_{\xi_0}^{\xi } (\tan \sigma_n - \tan (\sigma_{n-1}+a)) d\xi} {\varkappa_n - \varkappa_{n-1}-a} > - (\xi - \xi_0)S\delta.
\end{equation}

From \eqref{a3} and \eqref{a4},
\begin{equation}\label{a5}
\frac{S \int_{\xi_0}^{\xi } (\tan \sigma_n - \tan (\sigma_{n-1}+a)) d\xi} {\varkappa_n - \varkappa_{n-1}-a}\ge -S\delta(\xi-\xi_0) > - 1 - e^{\frac{1}{p}(\Delta_n-\Delta_{n+1}-\omega  )}. 
\end{equation}

However,
\[
\mathcal{G}K\varkappa_n(\xi) - \mathcal{G}K\varkappa_{n-1}(\xi)
= S \int_{\xi_0}^{\xi } (\tan \sigma_n - \tan \sigma_{n-1}) d\xi \ge S \int_{\xi_0}^{\xi } (\tan \sigma_n - \tan (\sigma_{n-1}+a)) d\xi.
\]

Therefore, \eqref{a5} becomes
\[
\frac{\mathcal{G}K\varkappa_n(\xi) - \mathcal{G}K\varkappa_{n-1}(\xi)}{\varkappa_n - \varkappa_{n-1}-a} >  - 1 - e^{\frac{1}{p}(\Delta_n-\Delta_{n+1}-\omega  )}.
\]

From physical considerations, $\varkappa(\xi)$ is inversely proportional to $\sigma(\xi)$. Hence, as per the assumption $(F_2)$,
\[
\dfrac{\mathcal{G}K\varkappa_{n}(\xi) - \mathcal{G}K\varkappa_{n-1}(\xi)}{\varkappa_n - \varkappa_{n-1}-a} < -1 + e^{\frac{1}{p}(\Delta_n-\Delta_{n+1}-\omega  )}.
\]

Thus from the above two inequalities,
\[
-1 - e^{\frac{1}{p}(\Delta_n-\Delta_{n+1}-\omega  )} < \frac{\mathcal{G}K\varkappa_n(\xi) - \mathcal{G}K\varkappa_{n-1}(\xi)}{\varkappa_n - \varkappa_{n-1}-a} < -1 + e^{\frac{1}{p}(\Delta_n-\Delta_{n+1}-\omega  )}.
\]

which yields that
\[
\left| \frac{\mathcal{G}K\varkappa_n(\xi) - \mathcal{G}K\varkappa_{n-1}(\xi)}{\varkappa_n - \varkappa_{n-1}-a} + 1 \right| < e^{\frac{1}{p}(\Delta_n-\Delta_{n+1}-\omega  )}
\quad \text{for all } \xi.
\]
Raising both sides to the power $p\ge1$, we have
\[
\left| \frac{\mathcal{G}K\varkappa_n(\xi) - \mathcal{G}K\varkappa_{n-1}(\xi)}{\varkappa_n - \varkappa_{n-1}-a} + 1 \right|^p < e^{(\Delta_n-\Delta_{n+1}-\omega  )}
\quad \text{for all } \xi.
\]
Thus, from Eq. \eqref{a1},
\[
\Delta_{n+1}+a^p < e^{(\Delta_n-\Delta_{n+1}-\omega  )} \Delta_n+a^p.
\]
\[\Rightarrow \max_{\xi }\left(\Delta_{n+1}+a^p\right)e^{\max_{\xi }(\Delta_{n+1}-\Delta_n)} < e^{-\omega  } \max_{\xi }\Delta_n+a^p.\]
\[\Rightarrow \frac{\max_{\xi }\Delta_{n+1}+a^p}{\max_{\xi }\Delta_n+a^p}\cdot\frac{e^{\max_{\xi }(\Delta_{n+1})}}{e^{\max_{\xi }\Delta_n}} < e^{-\omega  } .\]

By Lemma \ref{remark1}, this satisfies the conditions of Theorem \ref{sf theorem}, when $F(\xi) = \log \xi+\xi$.


Hence \eqref{aa2} converges if $(F_1)$ and $(F_2)$ are satisfied.
\end{proof}

\end{document}
